% SGA 3, Expose III, Loop 1 English reconstruction.
% Scope: Corollary 0.7 bis through Corollary 0.10 and its transition,
% stopping immediately before Proposition 0.11.
% Authority: Polo--Gille Exp3-14oct24.pdf, component-local pages 17--19
% (combined-reader pages 127--129; printed pages 117--119). Force-OCR is
% a drafting locator; English wording was compared with Jacob Reinhold's
% CC BY 4.0 Markdown lineage at commit
% e7a259f3f8608ad3edf9bf6eead3fd504dd2d23e. The PDF controls prose,
% mathematics, numbering, notes, and diagram images.

\setcounter{footnote}{26}

% Source page: Exp3 component-local p. 17; combined-reader p. 127; printed p. 117.
\medskip
\noindent\SGAInlineRefTarget{sga3:iii:corollary:0-7-bis}\textbf{推论 0.7 bis.}
\textit{设 \(X\) 是一个 \(S\)-概形, 使得 \(X_{\mathscr J}\) 是一个
\(S_{\mathscr J}\)-幺半群. 则 \(L_X\) 具有一个
\(S\)-幺半群结构, 且有 \(S\)-幺半群的分裂正合列:}

\SGAThreeNativeDiagram[0.50\textwidth]
  {figures/exp3/diagram_024_exp3_p017_corollary07bis_split_sequence.tex}
  {\SGARefLink{sga3:iii:corollary:0-7-bis}{推论 0.7 bis} 中 \(S\)-幺半群的分裂正合列.}
  {fig:exp3-corollary-07bis-split-sequence}
  {Polo--Gille Exp3-14oct24.pdf, 分篇页 17; 合订本页 127; 印刷页 117}

\textit{在 \(L'_X\) 上诱导的幺半群运算与其阿贝尔群结构一致.
特别地, 若 \(X_{\mathscr J}\) 是一个 \(S_{\mathscr J}\)-群,
则 \(L_X\) 是一个 \(S\)-群, 且为 \(X^+\) 与 \(L'_X\) 的半直积.}

\noindent\textit{证明.}
由于 \(X_{\mathscr J}\) 是一个 \(S_{\mathscr J}\)-幺半群,
\[
 X^+=\prod_{S_{\mathscr J}/S}X_{\mathscr J}
\]
是一个 \(S\)-幺半群; 事实上, 对每个 \(Y\to S\), 有
\(X^+(Y)=X_{\mathscr J}(Y_{\mathscr J})\).
对每个 \(S\)-概形 \(Y\), 以 \(\widetilde Y_{\mathscr J}\)
表示对应于拟凝聚 \(\mathscr O_{Y_{\mathscr J}}\)-代数
\[
 \mathscr O_{Y_{\mathscr J}}\oplus\mathscr J\mathscr O_Y,
\]
的仿射 \(Y_{\mathscr J}\)-概形; 此代数即与滤过
\(\mathscr O_Y\supset\mathscr J\mathscr O_Y\) 关联的分次代数.
于是 \(L_X(Y)\) 可与 \(X_{\mathscr J}(\widetilde Y_{\mathscr J})\)
等同, 而 \(L'_X(Y)\) 可与由\emph{零截面}
\(Y_{\mathscr J}\to\widetilde Y_{\mathscr J}\) 诱导的态射
\[
 p:X_{\mathscr J}(\widetilde Y_{\mathscr J})
   \longrightarrow X_{\mathscr J}(Y_{\mathscr J})
\]
的核等同. 等价地, 该截面由在理想 \(\mathscr J\mathscr O_Y\)
上为零的 \(\mathscr O_{Y_{\mathscr J}}\)-代数态射
\[
 \mathscr O_{\widetilde Y_{\mathscr J}}
   \longrightarrow\mathscr O_{Y_{\mathscr J}}
\]
诱导. 因此, 对每个 \(Y\to S\), 有如下关于 \(Y\) 具函子性的
幺半群分裂正合列:

\SGAThreeNativeDiagram[0.60\textwidth]
  {figures/exp3/diagram_025_exp3_p017_corollary07bis_pointwise_split_sequence.tex}
  {\(L_X(Y)\) 的逐点分裂正合列.}
  {fig:exp3-corollary-07bis-pointwise-sequence}
  {Polo--Gille Exp3-14oct24.pdf, 分篇页 17; 合订本页 127; 印刷页 117}

只须再证明, \(L'_X\) 上诱导的幺半群运算与其阿贝尔群结构一致.
设 \(\mu\) 是 \(L_X\) 的幺半群运算, \(e\) 是其单位截面.
必须证明, 对每个 \(m,m'\in L'_X(Y)\), 有
\[
 \mu(m\mathbin{\cdot}e,m'\mathbin{\cdot}e)
   =(m+m')\mathbin{\cdot}e.
\]

这有两种证明方法. 第一种方法是重复
（\SGARefLink{sga3:iii:corollary:0-1-10}{0.1.10} \((\ast)\)）的证明,
将其中的态射 \(f:X\to W\) 替换为
\[
 \mu:L_X\times_S L_X\longrightarrow L_X.
\]
将 \(X^+(Y)=X_{\mathscr J}(Y_{\mathscr J})\) 与其在 \(s\) 下于
\(L_X(Y)=X_{\mathscr J}(\widetilde Y_{\mathscr J})\) 中的像等同.
于是对所有 \(g,g'\in X_{\mathscr J}(Y_{\mathscr J})\) 和
\(m,m'\in L'_X(Y)\), 得
\begin{equation}
 \mu(m\mathbin{\cdot}g,m'\mathbin{\cdot}g')
 =L_\mu^{(g,g')}(m,m')\mathbin{\cdot}\mu(g,g').
 \tag{$\star$}\label{sga3:iii:corollary:0-7-bis:equation:star}
\end{equation}
这里 \(L_\mu^{(g,g')}\) 表示 \(\mu\) 在点 \((g,g')\) 处的导出态射;
换言之, \(\widetilde Y_{\mathscr J}\) 通过 \((g,g')\)
位于 \(L_X\times_S L_X\) 上. 特别地,
\[
 \mu(m\mathbin{\cdot}e,m'\mathbin{\cdot}e)
 =L_\mu^{(e,e)}(m,m')\mathbin{\cdot}e.
\]
而
\[
 L_\mu^{(e,e)}(m,m')=L_{\ell_e}(m')+L_{r_e}(m),
\]
其中 \(\ell_e\) 和 \(r_e\) 分别表示由 \(e\) 给出的左平移和右平移.
它们都是 \(X_{\mathscr J}\) 的恒等映射, 因而
\(L_\mu^{(e,e)}(m,m')=m+m'\).

另一种方法是如下论证（参见 [DG70], \S II.4,
\SGARefExternal{sga3:iii:theorem:3-5}{定理~3.5} 的证明）.
由\SGARefLink{sga3:iii:lemma:0-1-11}{引理~0.1.11},
\(X^+\) 与 \(L_X\) 的构造都与积可交换, 所以 \(L'_X\) 的构造也如此.
由此可知, \(\mu\) 诱导的态射
\[
 \mu':L'_X\times L'_X\longrightarrow L'_X
\]
是关于阿贝尔群结构的同态. 此时, 第~II 篇的
\SGARefLink{sga3:ii:lemma:3-10}{引理~3.10} 表明,
\(\mu'\) 与阿贝尔群运算一致.

% Source page: Exp3 component-local p. 18; combined-reader p. 128; printed p. 118.
\medskip
\noindent\SGAInlineRefTarget{sga3:iii:paragraph:0-8}\textbf{0.8.}
\footnote{编者注: 在此插入编号 0.8, 以标明从这里回到原文.}
现设 \(X\) 是一个 \(S\)-概形, 使得 \(X_{\mathscr J}\) 是一个
\(S_{\mathscr J}\)-群. 假设存在一个 \(S\)-态射
\[
 P:X\times_S X\longrightarrow X
\]
其基变换
\[
 P_{\mathscr J}:
 X_{\mathscr J}\times_{S_{\mathscr J}}X_{\mathscr J}
 \longrightarrow X_{\mathscr J}
\]
是 \(X_{\mathscr J}\) 的群运算. 一个重要特例是 \(X\) 为
\(S\)-群且 \(P\) 为其群运算的情形. 我们得到态射
\[
 L_P:L_X\times_S L_X\simeq L_{X\times_S X}\longrightarrow L_X.
\]
事实上, 这一态射不依赖于 \(P\), 因为它可以由
\(X_{\mathscr J}\) 的群运算 \(P_{\mathscr J}\) 计算出来,
下面即说明这一点.\footnote{编者注: 以下对原文作了展开.}

由\SGARefLink{sga3:iii:proposition:0-2:item:ii}{命题~0.2(ii)} 和 (iii),
对每个 \(Y\to S\), 每个 \(x,x'\in X(Y)\) 以及每个
\(m,m'\in L'_X(Y)\), 有
\[
\begin{aligned}
 P(m\mathbin{\cdot}x,m'\mathbin{\cdot}x')
 &=P\bigl((m,m')\mathbin{\cdot}(x,x')\bigr)\\
 &=L_P^{(x,x')}(m,m')\mathbin{\cdot}\mu(g,g'),
\end{aligned}
\]
其中 \(g\) 和 \(g'\) 分别为 \(x\) 和 \(x'\) 在 \(X^+(Y)\) 中的像.
此外（参见 \SGARefLink{sga3:iii:corollary:0-10}{0.10} 的证明）,
\(L_P^{(x,x')}=L_\mu^{(g,g')}\), 而由
\SGARefLink{sga3:iii:corollary:0-7-bis}{0.7 bis} \((\star)\),
它是在 \(L_X(Y)\) 中由等式
\[
 L_\mu^{(g,g')}(m,m')\mathbin{\cdot}\mu(g,g')
   =\mu(m\mathbin{\cdot}g,m'\mathbin{\cdot}g')
\]
定义的 \(L'_X(Y)\) 的元素.

以 \(\times\) 代替 \(\mu\) 表示 \(L_X\) 的群运算,
并以 \(\operatorname{Ad}\) 表示 \(X^+\) 在 \(L'_X\) 上的\emph{伴随作用}.
此作用经由 \(X_0\) 分解, 且由 \(X_0\) 在
\(\omega^1_{X_0/S_0}\) 上的伴随作用诱导. 则
\[
\begin{aligned}
 L_\mu^{(g,g')}(m,m')\times g\times g'
 &=m\times g\times m'\times g'\\
 &=\bigl(m\times\operatorname{Ad}(g)(m')\bigr)\times g\times g',
\end{aligned}
\]
因而
\[
 L_P^{(x,x')}(m,m')=m\times\operatorname{Ad}(g)(m').
\]
我们由此证明了:

\medskip
\noindent\SGAInlineRefTarget{sga3:iii:proposition:0-8}\textbf{命题 0.8.}
\textit{设 \(P:X\times_S X\to X\) 是一个 \(S\)-态射,
使得 \(P_{\mathscr J}\) 赋予 \(X_{\mathscr J}\) 一个
\(S_{\mathscr J}\)-群结构. 以 \(\times\) 表示 \(L'_X\) 的群运算,
以 \((m,x)\mapsto m\mathbin{\cdot}x\) 表示定义 \(L'_X\)
在 \(X\) 上的作用的态射 \(L'_X\times_S X\to X\), 并设}
\[
 \operatorname{Ad}:X^+\longrightarrow
 \operatorname{Aut}_{S\text{-}\mathrm{gr}}(L'_X)
\]
\textit{是 \(X^+\) 在 \(L'_X\) 上的伴随作用,
它由 \(X_0\) 在 \(\omega^1_{X_0/S_0}\) 上的伴随作用诱导.
则对每个 \(S'\to S\), 每个 \(x,x'\in X(S')\) 和每个
\(m,m'\in L'_X(S')\), 有}
\begin{equation}
 P(m\mathbin{\cdot}x,m'\mathbin{\cdot}x')
 =\bigl(m\times\operatorname{Ad}(p_X(x))(m')\bigr)
    \mathbin{\cdot}P(x,x').
 \tag{0.8.1}\label{sga3:iii:proposition:0-8:equation:0-8-1}
\end{equation}

\textit{若 \(X\) 是一个 \(S\)-群, 我们将以 \(\ast\) 表示其运算,
以 \(e\) 表示其单位截面, 以 \(i\) 表示由下式定义的 \(S\)-态射:}
\[
 i(m)=m\mathbin{\cdot}e
\]
\textit{此式对每个 \(S'\to S\) 及每个 \(m\in L'_X(S')\) 成立.}

% Source page: Exp3 component-local p. 19; combined-reader p. 129; printed p. 119.
\medskip
\noindent\SGAInlineRefTarget{sga3:iii:corollary:0-9}\textbf{推论 0.9.}
\textit{设 \(X\) 是一个 \(S\)-群. 则 \(X^+\) 自然成为一个
\(S\)-群, 且 \(p_X\) 是 \(S\)-群态射. 此外, \(S\)-态射}
\[
 i:L'_X\longrightarrow X,
 \qquad m\longmapsto m\mathbin{\cdot}e,
\]
\textit{是从 \(L'_X\) 到 \(\operatorname{Ker}(p_X)\) 的
\(S\)-群同构, 而且对每个 \(S'\to S\), 每个
\(x'\in X(S')\) 以及每个 \(m\in L'_X(S')\), 有}
\begin{equation}
 m\mathbin{\cdot}x'
 =(m\mathbin{\cdot}e)\ast x'
 =i(m)\ast x'.
 \tag{0.9.1}\label{sga3:iii:corollary:0-9:equation:0-9-1}
\end{equation}

前两个断言已经在 \SGARefLink{sga3:iii:remark:0-1-2}{0.1.2} 中证明.
由于 \(X\) 是 \(X^+\) 上关于
\(L_X=L'_X\times_S X^+\) 的形式主齐性空间,
态射 \(i\) 确实是从 \(L'_X\) 到 \(p_X\) 的核的 \(S\)-函子同构.
\(i\) 是群态射以及公式（\SGARefLink{sga3:iii:corollary:0-9:equation:0-9-1}{0.9.1}）,
分别由（\SGARefLink{sga3:iii:proposition:0-8:equation:0-8-1}{0.8.1}）取
\(x=x'=e\) 及取 \(x=e,\ m'=1\) 得到.

\medskip
\noindent\SGAInlineRefTarget{sga3:iii:corollary:0-10}\textbf{推论 0.10.}
\textit{设 \(X\) 是一个 \(S\)-群. 沿用上述记号,
对每个 \(S'\to S\), 每个 \(x\in X(S')\) 及每个
\(m'\in L'_X(S')\), 有}
\begin{equation}
 x\ast i(m')\ast x^{-1}
 =i\bigl(\operatorname{Ad}(p_X(x))(m')\bigr).
 \tag{0.10.1}\label{sga3:iii:corollary:0-10:equation:0-10-1}
\end{equation}
此式由等式 \(i(m')\ast x^{-1}=m'\mathbin{\cdot}x^{-1}\)
及（\SGARefLink{sga3:iii:proposition:0-8:equation:0-8-1}{0.8.1}）取
\(m=1\) 和 \(x'=x^{-1}\) 得到.

因此, 当 \(X\) 为 \(S\)-群时, 我们已明确确定了 \(X\to X^+\)
的核以及 \(X\) 的内自同构在此核上的作用.
现在将看到, 对某些未必可表的 \(S\)-群函子, 也能做同样的事.
有一种情形将对我们有用, 即\SGARefLink{sga3:i:subsection:1:7}{第~I 篇, 1.7}
中的函子 \(\underline{\operatorname{Aut}}\). 我们立刻在下一命题中陈述这一结果.

% Continue with Proposition 0.11.
